\chapter{Two qubit system: looking for entanglement}

\section{Convention, states and master equation}

We consider a system of $N$ qubits (with particular emphasis on the $N=2$ case), each described in the computational basis
$\{\ket{0},\ket{1}\}$.
Throughout this work we adopt the convention
\begin{equation}
	\sigma_z \ket{1} = + \ket{1}, 
	\qquad
	\sigma_z \ket{0} = - \ket{0},
\end{equation}
so that the excited state $\ket{1}$ corresponds to the positive eigenvalue of $\sigma_z$.
This choice fixes all relative signs in the collective spin operators and in the dissipative dynamics.

The single-qubit spin operators are defined as
\begin{equation}
	S_\alpha^{(i)} = \frac{1}{2}\,\sigma_\alpha^{(i)},
	\qquad \alpha \in \{x,y,z\},
\end{equation}
where the superscript $i$ labels the qubit.
Collective spin operators are introduced as
\begin{equation}
	J_\alpha = \sum_{i=1}^{N} S_\alpha^{(i)}.
\end{equation}

When needed, spin coherent states (CSS) are written as
\begin{equation}
	\ket{\theta,\phi}
	=
	\bigotimes_{i=1}^{N}
	\left(
	e^{i\phi}\sin\frac{\theta}{2}\,\ket{1}_i
	+
	\cos\frac{\theta}{2}\,\ket{0}_i
	\right),
\end{equation}
which corresponds to a product state with all spins pointing along the direction
$\hat{n} = (\sin\theta\cos\phi,\sin\theta\sin\phi,\cos\theta)$ on the Bloch sphere.
In particular, the eigenstates of $\sigma_x$ are
\begin{equation}
	\ket{+}=\frac{\ket{0}+\ket{1}}{\sqrt{2}},
	\qquad
	\ket{-}=\frac{\ket{0}-\ket{1}}{\sqrt{2}}.
\end{equation}

	We use the standard Bell states
\begin{align}
	\ket{\phi_\pm}&=\frac{\ket{00}\pm\ket{11}}{\sqrt{2}}
	\quad(\text{even $z$-parity}),\\
	\ket{\psi_\pm}&=\frac{\ket{01}\pm\ket{10}}{\sqrt{2}}
	\quad(\text{odd $z$-parity}).
\end{align}

The dynamics of the system is described within the Markovian open-quantum-system framework.
Under the standard assumptions of weak system--environment coupling, short bath correlation time
(Born--Markov approximation), and complete positivity, the reduced density operator $\rho(t)$
obeys a Gorini--Kossakowski--Sudarshan--Lindblad (GKSL) master equation,
\begin{equation}
	\frac{d\rho}{dt}
	= - i \left[ H, \rho \right] + \sum_k \mathcal{D}[L_k]\rho^{(c)}\,dt = 
	- i \left[ H, \rho \right]
	+
	\sum_k
	\left(
	L_k \rho L_k^\dagger
	-
	\frac{1}{2}
	\left\{
	L_k^\dagger L_k,\rho
	\right\}
	\right).
	\label{eq:Lindblad}
\end{equation}
where $\mathcal{D}[A]\rho = A\rho A^\dagger - \tfrac{1}{2}\{A^\dagger A,\rho\}$. Here $H$ is the system Hamiltonian,
and $\{L_k\}$ are the Lindblad or collapse operators describing dissipative processes and continuous measurements.

Equation~\eqref{eq:Lindblad} generates a completely positive and trace-preserving (CPTP) dynamical semigroup.

In this work we focus on dissipative channels associated with spin components along the $z$ direction.
For two qubits, a central role is played by collective jump operators of the form
\begin{equation}
	\hat{L}_\pm
	=
	\sqrt{\gamma}\,
	\frac{\sigma_z^{(1)} \pm \sigma_z^{(2)}}{\sqrt{2}}
	\label{eq:Lindblad_operators}
\end{equation}

where $\gamma$ sets the overall dissipation or measurement rate.

These operators naturally arise in measurement schemes where emissions from the two qubits are interfered
(e.g.\ via a beam splitter) before detection, such that the measurement does not reveal which-path information.
The two channels correspond to symmetric and antisymmetric combinations of local $z$-spin operators.
While the master equation~\eqref{eq:Lindblad} uniquely determines the evolution of the density operator,
its stochastic unravelling into quantum trajectories is not unique (see Appendix A)
Different measurement schemes (photodetection, homodyne, heterodyne) correspond to different
stochastic evolutions at the level of pure states, all reproducing the same ensemble-averaged dynamics.

In the photodetection unravelling associated with the jump operators~\eqref{eq:Lindblad_operators},
the system undergoes non-unitary no-click evolution interrupted by quantum jumps.
Each jump corresponds to the application of $L_\pm$ to the state vector,
followed by normalization, and reveals information about collective spin observables and parity.

This measurement backaction is the key mechanism exploited in this work for the generation
and stabilization of entanglement and for the study of metrological quantities such as
concurrence and spin squeezing.



\section{Photodetection}
\label{sec:photodetection}

A Markovian master equation admits a stochastic \emph{unravelling} in terms of
measurement-conditioned quantum trajectories. In the case of photodetection
(quantum-jump monitoring), the measurement record is a set of counting processes
$\{N_k(t)\}_k$ associated with the monitored output channels.
In an infinitesimal time step $dt$, the increments
\begin{equation}
	dN_k(t) \in \{0,1\}
\end{equation}
represent, respectively, a \emph{no-click} event or a \emph{click} in detector $k$.
They are Poissonian random variables with conditional mean
\begin{equation}
	\mathbb{E}\!\left[dN_k(t)\right]
	=
	\eta\,\mathrm{Tr}\!\left[L_k^\dagger L_k\,\rho^{(c)}(t)\right]dt ,
	\label{eq:poisson_mean_general}
\end{equation}
where $\eta\in[0,1]$ is the detection efficiency and $\rho^{(c)}(t)$ denotes the
conditional (normalized) system state.

For jump monitoring, the conditional evolution can be written as a piecewise
deterministic process: a smooth non-unitary evolution conditioned on $dN_k=0$,
interrupted by stochastic jumps when $dN_k=1$.
A convenient It\^o form of the stochastic master equation (SME) is
\begin{equation}
	\begin{aligned}
		d\rho^{(c)}
		&=
		-i[H,\rho^{(c)}]\,dt
		+ (1-\eta)\sum_k \mathcal{D}[L_k]\rho^{(c)}\,dt \\
		&\quad
		+ \sum_k
		\left(
		\frac{L_k\rho^{(c)}L_k^\dagger}{\mathrm{Tr}[L_k^\dagger L_k\rho^{(c)}]}
		-
		\rho^{(c)}
		\right)dN_k(t),
	\end{aligned}
	\label{eq:SME_jump_general}
\end{equation}

By averaging Eq.~\eqref{eq:SME_jump_general} over the measurement outcomes one
recovers the unconditional GKSL master equation.


The resulting photodetection SME reads
\begin{equation}
	d\varrho_t^{(c)}
	=
	-i\omega[J_z,\varrho_t^{(c)}]dt
	+ (1-\eta)\frac{\gamma}{2}\sum_j \mathcal{D}[\sigma_z^{(j)}]\varrho_t^{(c)}dt
	+ \sum_j\left(\sigma_z^{(j)}\varrho_t^{(c)}\sigma_z^{(j)}-\varrho_t^{(c)}\right)dN_j,
	\label{eq:AG_PD_SME}
\end{equation}
with $\mathbb{E}[dN_j]=\eta\gamma\,dt$.
A crucial remark (highly relevant for interpreting measurement records) is that, \textbf{\textit{since
$\sigma_z^{(j)}$ is unitary, the jump statistics is \emph{state-independent}: the records
$N_j(t)$ contain ``only noise'' and no information about the system state, and the SME is
linear in $\varrho_t^{(c)}$.}}
A Kraus-operator viewpoint over a time step $dt$ makes the structure transparent:
\begin{equation}
	M_0 = \mathbb{I} - iH_{\mathrm{eff}}dt,
	\qquad
	M_k = \sqrt{\eta\,dt}\,L_k,
\end{equation}
with non-Hermitian effective Hamiltonian
\begin{equation}
	H_{\mathrm{eff}} = H - \frac{i}{2}\eta\sum_k L_k^\dagger L_k,
\end{equation}
while the unmonitored fraction $(1-\eta)$ contributes the deterministic term in Eq.~\eqref{eq:SME_jump_general}.


\subsection{First jump from \texorpdfstring{$\ket{++}$}{|++>}: projection to definite parity subspace}

Since \(\sigma_z^{(1)}\) and \(\sigma_z^{(2)}\) commute and \((\sigma_z^{(i)})^2=\mathbb{I}\),
\begin{align}
	\big[\sigma_z^{(1)},\Pi_z\big]
	&=\sigma_z^{(1)}\sigma_z^{(1)}\sigma_z^{(2)}-\sigma_z^{(1)}\sigma_z^{(1)}\sigma_z^{(2)}=0,\\
	\big[\sigma_z^{(2)},\Pi_z\big]
	&=\sigma_z^{(2)}\sigma_z^{(1)}\sigma_z^{(2)}-\sigma_z^{(1)}\sigma_z^{(2)}\sigma_z^{(2)}
	=\sigma_z^{(1)}-\sigma_z^{(1)}=0.
\end{align}
Therefore \([\sigma_z^{(1)}\pm \sigma_z^{(2)},\Pi_z]=0\), and by \eqref{eq:Lindblad_operators} we have
\begin{equation}
	[\hat{C}_\pm,\Pi_z]=0.
\end{equation}
A jump produced by \(\hat{C}_\pm\) preserves the $z$-parity: if a pre-jump pure state \(\ket{\psi}\) is a parity eigenstate, so is the post-jump normalized state : 
\(\ket{\psi'}=\hat{C}_\pm\ket{\psi}/\norm{\hat{C}_\pm\ket{\psi}}\).
For pure-state jumps we use the standard update rule
\begin{equation}
	%\ket{\psi'}=\frac{C_k\ket{\psi}}{\sqrt{\matrixel{\psi}{C^{\dagger} }_k {\psi}}}\,,
	\qquad p_k=\matrixel{\psi}{C_k^\dagger C_k}{\psi}\,\mathrm{d}t,
	\label{eq:jump_rule}
\end{equation}
with \(k\in\{+,-\}\) and \(\mathrm{d}t\) a short time step.

It is convenient to evaluate the \emph{direction} of the unnormalized state \(\hat{C}_\pm\ket{++}\) using the numerators of \eqref{Cpm}:
\begin{align}
	(\sigma_z^{(1)}+\sigma_z^{(2)})\ket{++}
	&=2\big(\ket{11}-\ket{00}\big)
	=-2\sqrt{2}\,\ket{\phi_-},\label{eq:sigplus_on_pp}\\
	(\sigma_z^{(1)}-\sigma_z^{(2)})\ket{++}
	&=2\big(\ket{10}-\ket{01}\big)
	=-2\sqrt{2}\,\ket{\psi_-}.
	\label{eq:sigminus_on_pp}
\end{align}
Hence the \emph{normalized} post-jump states are
\begin{equation}
	\ket{++}\xrightarrow{C_+}-\ket{\phi_-},
	\qquad
	\ket{++}\xrightarrow{C_-}-\ket{\psi_-}.
	\label{eq:first_step}
\end{equation}
Overall minus signs are global phases and carry no physical effect; we keep them to make the subsequent alternation visually clear.

Equations \eqref{eq:first_step} show that the first click already fixes the $z$-parity: $C_+$ lands in the even sector (Bell-$\phi$ states), while $C_-$ lands in the odd sector (Bell-$\psi$ states).
From \eqref{Cpm},
\begin{equation}
	C_\pm^\dagger C_\pm
	=\frac{\gamma}{2}\big(S_z^{(1)}\pm S_z^{(2)}\big)^2
	=\frac{\gamma}{2}\left[\big(S_z^{(1)}\big)^2+\big(S_z^{(2)}\big)^2\pm2S_z^{(1)}S_z^{(2)}\right].
\end{equation}
Using \(S_z^2=\tfrac{1}{4}\,\mathbb{I}\) and \(S_z^{(1)}S_z^{(2)}=\tfrac{1}{4}\,\sigma_z^{(1)}\sigma_z^{(2)}=\tfrac{1}{4}\,\Pi_z\), we obtain the identity:
\begin{equation}
	C_\pm^\dagger C_\pm
	=\frac{\gamma}{4}\left(\mathbb{I}\pm \Pi_z\right).
	\label{eq:rate_identity}
\end{equation}
Therefore, for any pure state \(\ket{\psi}\), the transition probability rate is:
\begin{equation}
	p_\pm(\ket{\psi})=\matrixel{\psi}{C_\pm^\dagger C_\pm}{\psi}\,\mathrm{d}t
	=\frac{\gamma}{4}\Big(1\pm \ev{\Pi_z}_{\psi}\Big)\,\mathrm{d}t.
	\label{eq:general_rate}
\end{equation}
For the state \(\ket{++}\), since \(\ev{\sigma_z}_{\ket{+}}=0\), we have \(\ev{\Pi_z}_{\ket{++}}=\ev{\sigma_z}_{\ket{+}}\,\ev{\sigma_z}_{\ket{+}}=0\), hence
\begin{equation}
	p_+(\ket{++})=p_-(\ket{++})=\frac{\gamma}{4}\,\mathrm{d}t.
	\label{eq:rates_pp}
\end{equation}
The two channels are equiprobable on the initial state, and the total jump rate is \(\tfrac{\gamma}{2}\,\mathrm{d}t\).

For any even-parity state (eigenvalue \(+1\) of \(\Pi_z\); in particular \(\ket{\phi_\pm}\)),
\begin{equation}
	p_+(\text{even})=\frac{\gamma}{2}\,\mathrm{d}t,\qquad p_-(\text{even})=0.
\end{equation}
For any odd-parity state (eigenvalue \(-1\); in particular \(\ket{\psi_\pm}\)),
\begin{equation}
	p_-(\text{odd})=\frac{\gamma}{2}\,\mathrm{d}t,\qquad p_+(\text{odd})=0.
\end{equation}

\subsection{Subsequent jumps: two-state cyclicity}
Using again only the numerators (because normalization is restored at each jump), we find the action on Bell states:
\begin{align}
	\hat{C}_+\ket{\phi_-}
	&=-\,\ket{\phi_+},&
	\hat{C}_+\ket{\phi_+}
	&=-\,\ket{\phi_-},\\[2mm]
	\hat{C}_-\ket{\psi_-}
	&=-\,\ket{\psi_+},&
	\hat{C}_-\ket{\psi_+}
	&=-\,\ket{\psi_-}.
\end{align}
After renormalization, this yields the deterministic cycles
\begin{equation}
	-\ket{\phi_-}\xrightarrow{C_+}\ket{\phi_+}\xrightarrow{C_+}-\ket{\phi_-}\xrightarrow{C_+}\cdots,
	\qquad
	-\ket{\psi_-}\xrightarrow{C_-}\ket{\psi_+}\xrightarrow{C_-}-\ket{\psi_-}\xrightarrow{C_-}\cdots.
	\label{eq:cycles}
\end{equation}
Because of $[C_\pm,\Pi_z]=0$, the ``wrong'' channel is forbidden inside each parity subspace.

\subsection{State by state transitions - properties}
Let us collect both the post-jump states and the \emph{unconditional} jump probabilities over a time step \(\mathrm{d}t\):
\begin{center}
	\renewcommand{\arraystretch}{1.3}
	\begin{tabular}{c|c|c}
		Pre-jump state & Channel & Post-jump state \& probability \\ \hline
		\(\ket{++}\) & \(C_+\) & \(-\ket{\phi_-}\), \(\;p_+=\dfrac{\gamma}{4}\mathrm{d}t\) \\
		\(\ket{++}\) & \(C_-\) & \(-\ket{\psi_-}\), \(\;p_-=\dfrac{\gamma}{4}\mathrm{d}t\) \\ \hline
		\(-\ket{\phi_-}\) & \(C_+\) & \(\ket{\phi_+}\), \(\;p_+=\dfrac{\gamma}{2}\mathrm{d}t\) \\
		\(\ket{\phi_+}\) & \(C_+\) & \(-\ket{\phi_-}\), \(\;p_+=\dfrac{\gamma}{2}\mathrm{d}t\) \\
		(even sector) & \(C_-\) & forbidden, \(\;p_-=0\) \\ \hline
		\(-\ket{\psi_-}\) & \(C_-\) & \(\ket{\psi_+}\), \(\;p_-=\dfrac{\gamma}{2}\mathrm{d}t\) \\
		\(\ket{\psi_+}\) & \(C_-\) & \(-\ket{\psi_-}\), \(\;p_-=\dfrac{\gamma}{2}\mathrm{d}t\) \\
		(odd sector) & \(C_+\) & forbidden, \(\;p_+=0\)
	\end{tabular}
\end{center}
The selection rule is a direct consequence of \([C_\pm,\Pi_z]=0\): jumps cannot mix parity sectors. Equation~\eqref{eq:rate_identity} makes this quantitative. Thus, once the first click has selected a parity subspace, only the corresponding channel can ever fire; conditioned on a click, the state update within that sector is deterministic as in \eqref{eq:cycles}.

Let's analyze the variances within each channel's cycle. For Bell states the correlation matrix is diagonal:
\[
\expval*{\hat\sigma_\alpha^{(1)}\!\otimes \hat\sigma_\alpha^{(2)}}=
\begin{cases}
	diag(+1, -1, +1) & \text{for }\ket{\phi_+}\ \text{with }\alpha=x,y,z,\\
	diag(-1, +1, +1) & \text{for }\ket{\phi_-},\\
	diag(+1, +1, -1) & \text{for }\ket{\psi_+},\\
	diag(-1, -1, -1) & \text{for }\ket{\psi_-}.
\end{cases}
\]
Using
\[
\hat J_\alpha^2=\tfrac{1}{2}\Big(\hat{\mathbb{I}}+\hat\sigma_\alpha^{(1)}\!\otimes \hat\sigma_\alpha^{(2)}\Big),
\qquad \expval*{\hat J_\alpha}=0\ \text{for Bell states},
\]
one gets
\[
\Delta J_\alpha^2=\expval*{\hat J_\alpha^2}
=\tfrac{1}{2}\Big(1+\expval*{\hat\sigma_\alpha^{(1)}\!\otimes \hat\sigma_\alpha^{(2)}}\Big).
\]
\medskip
So we summarize in a table of variances \(\Delta J_\alpha^2\), for the two cycles):
\[
\begin{array}{c|ccc}
	\text{state} & \Delta J_x^2 & \Delta J_y^2 & \Delta J_z^2\\
	\hline
	\ket{\phi_-} & 0 & 1 & 1\\
	\ket{\phi_+} & 1 & 0 & 1\\
	\hline
	\ket{\psi_-} & 0 & 0 & 0\\
	\ket{\psi_+} & 1 & 1 & 0
\end{array}
\]	


%%%%%%%%%%%%%%%%%%%%%%%%%%%%%%%%%%%%%%%%%%%%%%%%%%%%%%%%%%%%%%%%%%%%%%%%%%%%%%%%%%%%%%%%%%%%%%%%%%%%
%%%%%%%%%%%%%%%%%%%%%%%%%%%%%%%%%%%%%%%%%%%%%%%%%%%%%%%%%%%%%%%%%%%%%%%%%%%%%%%%%%%%%%%%%%%%%%%%%%%%



\subsection{Ensemble-averaged quantities over trajectories}

For a pure two-qubit state \(\ket{\psi}=a\ket{00}+b\ket{01}+c\ket{10}+d\ket{11}\),
the concurrence is \(\mathcal{C}(\ket{\psi})=2\,\lvert ad-bc\rvert\). For a separable state it is 0 and for any Bell states it is 1.
Hence \(\mathcal{C}(\ket{++})=0\) and \(\mathcal{C}(\ket{\phi_\pm})=\mathcal{C}(\ket{\psi_\pm})=1\).

Along each trajectory the concurrence is a step process: it remains \(0\) until the \emph{first} click, and jumps to \(1\) thereafter (all subsequent clicks keep the state maximally entangled within the same parity subspace).

The waiting time to the first click is exponential with rate
\[
\lambda=\Tr\!\Big[\big(\hat C_+^\dagger \hat C_+ + \hat C_-^\dagger \hat C_-\big)\, \cdot \rho_{++}\Big]=\frac{\gamma}{2},
\]
so the survival probability of no click up to time \(t\) is \(P_0(t)=e^{-\lambda t}=e^{-(\gamma/2)t}\).
Hence the ensemble-averaged concurrence over many trajectories is
\[
\boxed{\ \overline{\mathcal{C}}(t)= prob (\text{at least one click by }t)=1-e^{-(\gamma/2)t}\ } 
\]
With \(\gamma=1\), \(\overline{\mathcal{C}}(t)=1-e^{-t/2}\).
